\scanpage{183}
\mspage{II.31}
% archive sheet number: 182
% block: upper left, paired title strokes | drawing
\margin{\figure{\begin{tikzpicture}[x=1mm,y=1mm]
\draw (0,0)--(0,17);\draw (4,-2)--(4,11);
\end{tikzpicture}}}
% block: top, lines 1-3 | text
\uline{Appendice.}

6. \uline{Groupes et Algèbres de Lie \inserted{analytiques \unsure{ou} algébriques \deleted{de}} associés aux représentations $\ell$-adiques de groupes profinis}

% block: upper third, line 4 | text
On se fixe un nb premier $\ell$.

% block: upper third, lines 5-10 | text
6.1. Rappelons que sur un groupe topologique \unsure{admettant} une structure de groupe analytique sur $\mathbb Q_\ell$, cette dernière est unique (et si le groupe n'est pas discret, $\ell$ est uniquement déterminé); \deleted{\unsure{et tout}} \unsure{tout} \inserted{\unsure{tout}} hom. \uline{continu} de groupes analytiques sur $\mathbb Q_\ell$ est automatiquement continu. Enfin, un s-groupe \uline{fermé} d'un groupe analytique (sur $\mathbb Q_\ell$) est encore analytique.

% block: middle third, lines 11-16 | text
6.2. Soit \deleted{$\unsure{G}$} $\pi$ un groupe profini. Si $G=\pi/V$ et $G'=\pi/V'$ sont analytiques sur $\mathbb Q_\ell$, alors $G''=G/V\cap V'$ aussi, car il \deleted{\unsure{s'injecte}} est isomorphe à un sous-groupe fermé de $G\times G'$. Donc les \deleted{$G/V$} quotients \uline{analytiques} de $G$ forment un système projectif de groupes analytiques

% block: middle third, line 17 | display
\[
\deleted{\unsure{\mathscr G_\ell\pi}}\deleted{\unsure{\mathscr L_\ell(\pi)|\mathscr G_\ell(\pi)}}\inserted{\pi^{\mathrm{an}}}=(G_\alpha)
\]

% block: middle to lower third, lines 18-21 | text
qu'on appellera ``syst. projectif de Lie'' associé à $\pi$. \unsure{Noter|Notons} que si on considère $\pi$ comme \deleted{\unsure{pro-groupe}} \deleted{\unsure{profini}} syst. projectif de groupes \deleted{\unsure{analytiques}}\inserted{finis, considérés} comme groupes analytiques, on a un hom évident

% block: lower third, line 22 | display
\[
\deleted{\unsure{\mathscr L_\ell}}\pi^{\mathrm{an}}\mathbin{-}\pi.
\]

% block: lower third, lines 23-27 | text
6.3. Les représentations $\ell$-adiques de $\pi$ (sur $\mathbb Z_\ell$, ou sur $\mathbb Q_\ell$) ci-dessus s'interprètent alors comme représentations \deleted{\unsure{topologiques}}\inserted{\unsure{continus|continues} (analytiques)} de $\unsure{\mathscr L_\ell(\pi)|\mathscr G_\ell(\pi)}$, i.e. la catégorie des modules $\ell$-adiques sur $\pi$ est la \deleted{limite}\inserted{lim} des catégories analogues
